%==============================================================================
%  Academic Report: Valuation and Risk Management of Synthetic CDOs
%  Author : Nguyen Duc Khanh Trinh (TOM)
%  Based on the Python working paper "CDO_Pipeline_Pro.ipynb"
%  References: Mounfield (2009), Li (2000), Vasicek (1987), JP Morgan (2004)
%==============================================================================
\documentclass[11pt,a4paper]{article}

% -------------------------  packages  ----------------------------------------
\usepackage[utf8]{inputenc}
\usepackage[T1]{fontenc}
\usepackage[english]{babel}
\usepackage[margin=1in]{geometry}
\usepackage{amsmath,amssymb,amsthm,mathtools}
\usepackage{bm}
\usepackage{graphicx}
\usepackage{booktabs}
\usepackage{array}
\usepackage{tabularx}
\usepackage{multirow}
\usepackage{caption}
\usepackage{subcaption}
\usepackage{hyperref}
\usepackage{xcolor}
\usepackage{microtype}
\usepackage{enumitem}
\usepackage{siunitx}
\usepackage{fancyhdr}
\usepackage{setspace}
\usepackage{float}
\usepackage[numbers,sort&compress]{natbib}
\usepackage{titlesec}
\usepackage{algorithm}
\usepackage{algpseudocode}

\hypersetup{
    colorlinks=true,
    linkcolor=blue!55!black,
    citecolor=teal!55!black,
    urlcolor=blue!60!black,
}

% -------------------------  theorems -----------------------------------------
\newtheorem{definition}{Definition}[section]
\newtheorem{proposition}[definition]{Proposition}
\newtheorem{remark}[definition]{Remark}

% -------------------------  headers ------------------------------------------
\pagestyle{fancy}
\fancyhf{}
\rhead{\small Synthetic CDO --- Valuation \& Risk}
\lhead{\small Working paper}
\cfoot{\thepage}

% -------------------------  maths shortcuts ----------------------------------
\newcommand{\E}{\mathbb{E}}
\newcommand{\PP}{\mathbb{P}}
\newcommand{\QQ}{\mathbb{Q}}
\newcommand{\R}{\mathbb{R}}
\newcommand{\N}{\mathbb{N}}
\newcommand{\1}{\mathbf{1}}
\newcommand{\dd}{\mathrm{d}}
\newcommand{\Nor}{\mathcal{N}}
\DeclareMathOperator{\LGD}{LGD}
\DeclareMathOperator{\ETL}{ETL}
\DeclareMathOperator{\DF}{DF}
\DeclareMathOperator{\CS}{CS01}

% -------------------------  title --------------------------------------------
\title{\bfseries Valuation and Risk Management of Synthetic Collateralized Debt Obligations\\
       \large A Monte-Carlo and Semi-Analytic Pipeline under the One-Factor
       Gaussian Copula}
\author{Nguyen Duc Khanh Trinh\\
        \small \texttt{nguyenduckhanhtrinh@gmail.com}\\[2pt]
        \small ISFA --- Credit Risk Working Paper}
\date{Valuation date: 10~September~2021}

%==============================================================================
\begin{document}
\maketitle

\begin{abstract}
\noindent
This paper develops an end-to-end pricing and risk-management pipeline for
\emph{synthetic Collateralized Debt Obligations} (CDOs) on a reference pool
of 651 single-name CDS issuers. The framework follows the standard sell-side
workflow: \emph{term-sheet \(\rightarrow\) market data \(\rightarrow\)
marginal credit curves \(\rightarrow\) dependence model
\(\rightarrow\) tranche cash-flow engine \(\rightarrow\) calibration
\(\rightarrow\) sensitivities \(\rightarrow\) stress \(\rightarrow\) dashboard}.
Marginal default probabilities are bootstrapped from a term-structure of CDS
quotes using an ISDA-consistent piecewise-constant hazard model. Dependence
is modelled via the one-factor Gaussian copula of \citet{li2000,vasicek1987}.
Tranche expected losses are computed with a Monte-Carlo engine that uses
antithetic variance reduction and benchmarked against the Vasicek large
homogeneous pool (LHP) closed-form solution. A base-correlation curve is
calibrated by bisection and then interpolated to price bespoke tranches.
We provide a complete sensitivity analysis (tranche-level and name-level
CS01, Rho01, DV01, Jump-to-Default) and a deterministic stress matrix
covering spread, recovery, correlation and a combined crisis scenario.
A historical portfolio dashboard from 2015 to 2021 illustrates the behaviour
of each tranche through the COVID-19 credit shock. The complete results
show that the pipeline is internally consistent (ISDA repricing
error \(<10^{-16}\) on 300 test names), that the standard
capital-structure monotonicity
\(s^{\star}_{\text{equity}}=1\,952.6\) bp
\(>s^{\star}_{\text{mezz}}=503.3\) bp
\(>s^{\star}_{\text{jsr}}=202.7\) bp
\(>s^{\star}_{\text{senior}}=85.4\) bp
is preserved at \(\rho=0.30\) on the 10~September~2021 snapshot,
and that the signs of \(\partial s^{\star}/\partial\rho\) match
economic intuition (equity long correlation with
\(\mathrm{Rho01}_{\mathrm{eq}}=-37.5\), senior short correlation
with \(\mathrm{Rho01}_{\mathrm{sen}}=+3.83\)).\medskip

\noindent\textbf{Keywords:} Synthetic CDO, Gaussian Copula, Base Correlation,
Monte-Carlo, Expected Tranche Loss, Credit Risk, CS01.\\[2pt]
\textbf{JEL Classification:} G12, G13, G32, C15, C63.
\end{abstract}

\tableofcontents
\bigskip\hrule\bigskip

%==============================================================================
\section{Introduction}\label{sec:intro}

Synthetic CDOs are structured credit products whose cash flows depend on
the number of defaults that occur in a reference portfolio of corporate
obligors during the life of the contract.  Since their popularisation in
the early 2000s and the well-documented distress of the
2007--2009 crisis \citep{mounfield2009}, they have become a benchmark
laboratory for applying dependence modelling and tail-risk quantification
to credit. Standardised tranches on the CDX and iTraxx indices still trade
actively and the same mathematical toolbox is used today for bespoke
tranche opportunities and for regulatory capital computations.

The central pricing difficulty is that a tranche's present value is not
determined by individual default probabilities alone but by the
\emph{joint} loss distribution of the pool. A correct dependence structure
is therefore indispensable. The industry solution remains
the one-factor Gaussian copula of \citet{li2000,vasicek1987} with the
base-correlation extension of \citet{jpmorgan2004}.  Despite well-known
limitations --- most notably the correlation smile --- the model retains
academic and practical appeal thanks to its tractability.

\paragraph{Contributions.}
The present paper documents a full reproducible implementation and
provides four practical contributions.
First, we describe a complete Python pipeline (attached Jupyter notebook
\texttt{CDO\_Pipeline\_Pro.ipynb}) that has no dependency on
\texttt{scipy}. Second, we provide a benchmark between the Monte-Carlo
tranche pricer and the Vasicek LHP semi-analytic formula on a realistic
heterogeneous pool of 651 issuers, thereby quantifying the pedagogical
heterogeneity error of LHP. Third, we build a bespoke pricing surface
via base-correlation interpolation. Fourth, we produce a complete risk
dashboard that includes tranche-level Greeks, name-level CS01, a
deterministic stress matrix and a historical tranche ETL walk through
the COVID-19 credit shock.

\paragraph{Roadmap.}
Section~\ref{sec:framework} introduces the pricing framework.
Section~\ref{sec:curves} describes the marginal credit-curve bootstrap.
Section~\ref{sec:copula} formalises the Gaussian copula and the LHP
limit. Section~\ref{sec:mc} presents the Monte-Carlo engine and the
antithetic variance reduction. Section~\ref{sec:calib} discusses the
base-correlation calibration and the bespoke pricing surface.
Section~\ref{sec:risk} details the risk analytics and stress matrix.
Section~\ref{sec:dashboard} provides the historical dashboard.
Section~\ref{sec:conclusion} concludes.

%==============================================================================
\section{The Pricing Framework}\label{sec:framework}

\subsection{Term sheet and market inputs}
The deal is specified by a reference portfolio
\(\{i=1,\dots,N\}\) of CDS issuers with equal notional
\(N_{i}=1/N\), a set of tranche attachment/detachment points
\(\{[A_{j},D_{j}]\}_{j}\), a quarterly coupon schedule with maturity
\(T=5\) years and a convention on recovery
\(R\). The market inputs are the term structure of CDS spreads,
\(\{s_{i}^{T_{k}}\}_{k=1}^{K}\), and a risk-free discount curve
\(\DF(t)\). We use a flat continuously compounded risk-free rate
\(r=2\%\) so that \(\DF(t)=e^{-rt}\), and the ISDA-standard
\(R=40\%\) (LGD \(=60\%\)) for senior unsecured.

\subsection{Reference portfolio}\label{ssec:portfolio}
The portfolio is reconstructed from the time series
\texttt{cds.csv} containing single-name 1Y, 2Y, 3Y, 5Y, 7Y and 10Y CDS
spreads over 1747 daily observations (January~2015 to September~2021)
for 6610 ticker-tenor pairs. After filtering to the 5Y pillar we
obtain a reference pool of \(N=651\) issuers (equal-weighted). The
valuation snapshot is taken at the last available date,
\textbf{10~September~2021}.
Figure~\ref{fig:cds_ts} displays the aggregate dynamics of the pool and
its cross-sectional heterogeneity, while Table~\ref{tab:capital_struct}
summarises the standardised iTraxx-like capital structure used throughout.

\begin{figure}[H]
    \centering
    % TODO: insert pool-average 5Y CDS time series figure here
    % \includegraphics[width=0.8\textwidth]{fig_cds_timeseries.png}
    \fbox{\parbox[c][5cm][c]{0.75\textwidth}{\centering\itshape
        Figure placeholder --- CDS time series: pool-average 5Y spread
        (January 2015 -- December 2021) with the March~2020 COVID-19
        shock highlighted.}}
    \caption{Time series of the pool-average 5Y CDS spread, 2015--2021.}
    \label{fig:cds_ts}
\end{figure}

\begin{table}[H]
\centering
\caption{Standard iTraxx-like capital structure used in the pipeline.}
\label{tab:capital_struct}
\begin{tabular}{lcccl}
\toprule
\textbf{Tranche} & \textbf{Attach} $A$ & \textbf{Detach} $D$ &
\textbf{Width} & \textbf{Economic role}\\
\midrule
Equity          & 0\%  & 3\%  & 3\% & First loss; \emph{long} correlation\\
Mezzanine       & 3\%  & 7\%  & 4\% & Short gamma; default clustering\\
Junior Senior   & 7\%  & 10\% & 3\% & Dual protection; rho neutral\\
Senior          & 10\% & 15\% & 5\% & Tail protection; \emph{short} correlation\\
\bottomrule
\end{tabular}
\end{table}

\subsection{Tranche loss function}
Denote \(L^{\text{pool}}(t)=\sum_{i}\LGD_{i}\,\1_{\{\tau_{i}\le t\}}/N\)
the cumulative pool loss as a fraction of the notional. The
cash-flow entitlement of tranche \([A,D]\) is the capped-and-floored
pool loss
\begin{equation}
L^{[A,D]}(t) \;=\; \min\!\bigl(\max(L^{\text{pool}}(t)-A,\,0),\;D-A\bigr).
\label{eq:tranche_loss}
\end{equation}
The \emph{expected tranche loss} \(\ETL(t) = \E[L^{[A,D]}(t)]\) is
then the central quantity to compute.

\subsection{Premium, protection and the fair running spread}
A synthetic CDO tranche has two cash-flow legs. The \emph{protection
leg} pays the increment in tranche loss on the payment schedule
\(\{t_{0}=0,t_{1},\dots,t_{K}=T\}\); the \emph{premium leg} pays the
running spread \(s\) on the \emph{outstanding} tranche notional
\((D-A)-L^{[A,D]}(t)\). Denoting
\(\Delta t_{k}=t_{k}-t_{k-1}\), a standard trapezoidal approximation
yields \citep{mounfield2009}
\begin{align}
V_{\text{prot}}(s)
 &= \sum_{k=1}^{K} \DF(t_{k})\,\bigl(\ETL_{k}-\ETL_{k-1}\bigr),
 \label{eq:prot_leg}\\[2pt]
V_{\text{prem}}(s)
 &= s\sum_{k=1}^{K} \DF(t_{k})\,\Delta t_{k}\,
    \Bigl((D-A) - \tfrac{1}{2}\bigl(\ETL_{k}+\ETL_{k-1}\bigr)\Bigr).
 \label{eq:prem_leg}
\end{align}
The \emph{fair running spread} \(s^{\star}\) equates the two legs,
\begin{equation}
s^{\star} \;=\;
 \frac{\displaystyle\sum_{k=1}^{K}\DF(t_{k})\,(\ETL_{k}-\ETL_{k-1})}
      {\displaystyle\sum_{k=1}^{K}\DF(t_{k})\,\Delta t_{k}\,
      \bigl((D-A)-\tfrac{1}{2}(\ETL_{k}+\ETL_{k-1})\bigr)}.
\label{eq:fair_spread}
\end{equation}
Equation~\eqref{eq:fair_spread} is the \emph{master equation} of the
pipeline: once the curve \((\ETL_{k})\) is available under the chosen
model, all pricing, calibration and sensitivity computations reduce to
applications of this identity.

%==============================================================================
\section{Marginal Credit Curves}\label{sec:curves}

For each issuer we need a survival curve \(Q_{i}(t)\) under the
risk-neutral measure \(\QQ\).  We document two specifications.

\subsection{Flat hazard baseline (credit triangle)}
Assuming a constant intensity \(\lambda_{i}\) and taking the limit of
continuous coupon payments on the CDS protection leg yields Hull's
\emph{credit triangle}
\begin{equation}
\lambda_{i} \;\approx\; \frac{s_{i}^{5Y}}{\,1-R\,},
\qquad Q_{i}(t)=e^{-\lambda_{i}\,t}.
\label{eq:flat_hazard}
\end{equation}
This closed form is convenient for the historical dashboard but
abstracts from the term structure.

\subsection{ISDA-consistent piecewise-constant bootstrap}
The production specification allows \(\lambda_{i}(t)\) to be
piecewise constant on the tenor grid
\(0=T_{0}<T_{1}<\cdots<T_{K}\). The survival function is
\begin{equation}
Q_{i}(t)=\exp\!\Bigl(-\sum_{k}\lambda_{i}^{(k)}\Delta_{k}(t)\Bigr),
\quad \Delta_{k}(t)=\max(\min(T_{k},t)-T_{k-1},0).
\label{eq:survival_pwc}
\end{equation}
For each tenor \(T_{k}\) the hazard \(\lambda_{i}^{(k)}\) is determined
by requiring the market quote to be consistent with the CDS par identity
\begin{equation}
\underbrace{\int_{0}^{T_{k}}s\,\DF(u)\,Q(u)\dd u}_{\text{premium annuity}}
\;=\;
\underbrace{(1-R)\int_{0}^{T_{k}}\DF(u)\bigl(-\dd Q(u)\bigr)}_{\text{protection}} .
\label{eq:cds_par}
\end{equation}
The bootstrap proceeds sequentially: fix
\(\lambda^{(1)},\dots,\lambda^{(k-1)}\) and apply a one-dimensional
bisection on \(\lambda^{(k)}\) until \eqref{eq:cds_par} is satisfied;
this is the ISDA Standard Model mechanics.

\paragraph{Self-consistency.}
After the bootstrap completes we verify that repricing the same CDS
using the calibrated curve returns a par quote with NPV
\(<10^{-5}\).  On a random sub-sample of 300 names from the pool we
obtain a maximum absolute NPV of \(\mathbf{8.33\!\times\!10^{-17}}\)
and a mean of \(\mathbf{2.95\!\times\!10^{-17}}\), with \textbf{100\%}
of names passing the threshold. The bootstrap is therefore
ISDA-consistent up to machine precision. Pool-average 5Y PDs are
\(5.574\%\) for the flat-hazard baseline and \(5.717\%\) for the
piecewise-constant bootstrap, with cross-sectional range
\([0.859\%,\,47.778\%]\).
Figure~\ref{fig:hazard_ts} illustrates the resulting term-structure of
hazard rates grouped by 5Y-PD quintile; the expected shape --- flat for
strong credits and upward-sloping for distressed names --- is recovered.

\begin{figure}[H]
    \centering
    % TODO: insert hazard term-structure figure
    % \includegraphics[width=0.8\textwidth]{fig_hazard_term_structure.png}
    \fbox{\parbox[c][5cm][c]{0.75\textwidth}{\centering\itshape
        Figure placeholder --- Average bootstrapped hazard rate
        \(\lambda^{(k)}\) versus tenor \(T_{k}\), grouped by 5Y-PD
        quintile.  Strong credits are flat and low; distressed names
        show an upward slope consistent with credit-deterioration
        expectations.}}
    \caption{Piecewise-constant hazard term structure by rating bucket.}
    \label{fig:hazard_ts}
\end{figure}

%==============================================================================
\section{The One-Factor Gaussian Copula}\label{sec:copula}

\subsection{Model specification}
Following \citet{li2000} and the structural interpretation of
\citet{vasicek1987}, for each obligor we define a latent asset return
\begin{equation}
A_{i}=\sqrt{\rho}\,M+\sqrt{1-\rho}\,\varepsilon_{i},
\qquad M,\varepsilon_{i}\stackrel{\mathrm{iid}}{\sim}\Nor(0,1),
\label{eq:latent}
\end{equation}
with \(\rho\in[0,1]\) the asset correlation. Default occurs when
\(A_{i}<\Phi^{-1}\!\bigl(p_{i}(T)\bigr)\). Conditional on the common
factor \(M=m\), defaults become independent with
\begin{equation}
p_{i}(T\mid M=m)=\Phi\!\Biggl(\frac{\Phi^{-1}(p_{i}(T))-\sqrt{\rho}\,m}
{\sqrt{1-\rho}}\Biggr).
\label{eq:cond_pd}
\end{equation}
Economic intuition: adverse systemic factors (\(m\) small, negative)
push all conditional PDs up simultaneously, generating
\emph{default clustering}. The parameter \(\rho\) governs the severity
of such clustering and, consequently, the thickness of the right tail
of the loss distribution.  A representative illustration with
marginal PD \(p=5\%\) and \(\rho=0.30\) gives
\(p(T\mid M\!=\!-2)=25.57\%\),
\(p(T\mid M\!=\!0)=2.47\%\) and
\(p(T\mid M\!=\!+2)=0.053\%\), which makes the non-linear
amplification in the tail explicit.

\subsection{The Large Homogeneous Pool (LHP) limit}
Under \(p_{i}=p\) for all \(i\) and \(N\to\infty\), the law of large
numbers gives the Vasicek loss limit
\begin{equation}
L^{\text{pool}}(T)\;\xrightarrow{\;p\;}\;
\LGD\cdot\Phi\!\Biggl(\frac{\Phi^{-1}(p)-\sqrt{\rho}\,M}{\sqrt{1-\rho}}
\Biggr),
\end{equation}
from which the closed-form loss CDF follows:
\begin{equation}
F_{L}(\ell)
 =\Phi\!\Biggl(\frac{\sqrt{1-\rho}\,\Phi^{-1}(\ell/\LGD)-\Phi^{-1}(p)}
{\sqrt{\rho}}\Biggr).
\label{eq:vasicek_cdf}
\end{equation}
This is the cornerstone of the Basel II IRB risk-weight formula
\citep{vasicek1987,bcbs2006}.  In the pipeline, \eqref{eq:vasicek_cdf}
is integrated with a combined Gauss--Hermite/Gauss--Legendre quadrature
(64 outer and 96 inner nodes) to compute the LHP expected tranche loss
and used as an independent benchmark against the Monte-Carlo output
and for the historical dashboard.

\begin{figure}[H]
    \centering
    % TODO: insert LHP loss density / correlation comparison
    % \includegraphics[width=0.8\textwidth]{fig_lhp_correlation.png}
    \fbox{\parbox[c][5cm][c]{0.75\textwidth}{\centering\itshape
        Figure placeholder --- Pool-loss density under the LHP model
        for \(\rho\in\{0.10,\,0.25,\,0.40\}\). The right tail fattens
        substantially with \(\rho\); this is the mechanism behind
        senior-tranche sensitivity to correlation.}}
    \caption{Pool-loss density as a function of asset correlation
    \(\rho\).}
    \label{fig:lhp_rho}
\end{figure}

%==============================================================================
\section{Monte-Carlo Pricing Engine}\label{sec:mc}

\subsection{Default-time simulation}
For a given correlation \(\rho\) and calibrated marginal PDs
\(\{p_{i}(T)\}\), a Monte-Carlo path is generated by drawing
\(M\sim\Nor(0,1)\) and \(\varepsilon_{i}\sim\Nor(0,1)\), forming
\(Y_{i}\) according to \eqref{eq:latent} and mapping to the default
time
\begin{equation}
U_{i}=\Phi(Y_{i}),\qquad
\tau_{i}=-\log\bigl(1-U_{i}\bigr)\,\big/\,\lambda_{i},
\label{eq:default_time}
\end{equation}
under the marginal exponential calibration consistent with the flat
hazard (\S\ref{sec:curves}). The same \texttt{tau\_world} realisation is
re-used across tranches, so that Monte-Carlo noise is common and does
not contaminate the relative ordering of tranche spreads.

\subsection{ETL on the payment grid}
On the quarterly schedule \(\Delta t = 0.25\), \(T=5Y\),
\(K=20\) coupons, each path produces a loss trajectory
\(\bigl(L^{\text{pool}}(t_{k})\bigr)_{k=0}^{K}\), which is transformed
via \eqref{eq:tranche_loss} into a tranche-loss path. Averaging across
paths yields \((\ETL_{k})_{k=0}^{K}\). Figure~\ref{fig:etl_paths}
displays the cumulative ETL of the four tranches.

\begin{figure}[H]
    \centering
    % TODO: ETL layering by tranche
    % \includegraphics[width=0.8\textwidth]{fig_etl_layering.png}
    \fbox{\parbox[c][5cm][c]{0.75\textwidth}{\centering\itshape
        Figure placeholder --- Cumulative expected tranche loss
        \(\ETL_{k}\) on the quarterly grid, by tranche:
        equity, mezzanine, junior senior and senior.}}
    \caption{Capital-structure layering of expected tranche losses.}
    \label{fig:etl_paths}
\end{figure}

\subsection{Antithetic variance reduction}
To reduce variance, for each systemic draw \(M=m\) we add the
antithetic draw \(M=-m\). Let \(\sigma^{2}_{\text{crude}}\) and
\(\sigma^{2}_{\text{anti}}\) denote the respective variances of the
ETL estimator; the variance reduction factor
\begin{equation}
\mathrm{VRF}\;=\;\sigma^{2}_{\text{crude}}\big/\sigma^{2}_{\text{anti}}
\end{equation}
measured on the mezzanine tranche with \num{25000} paths is
\(\mathrm{VRF}\approx 0.98\).  For this pool the antithetic variate
therefore offers essentially no variance reduction: the mezzanine
loss profile is approximately symmetric in the systemic factor, so
\(L^{[A,D]}(M)+L^{[A,D]}(-M)\) is close to \(2\,\E[L^{[A,D]}]\) and
the antithetic pair contributes little beyond the crude estimator.
This is an important empirical finding: \emph{antithetic variance
reduction is not automatic} and the VRF should be measured on each
specific pool before being claimed as a speed-up.

\subsection{LHP vs Monte-Carlo benchmark}
On the heterogeneous 651-name pool, we compare the Monte-Carlo ETL
against the LHP formula evaluated at the pool-average PD
\begin{equation}
p^{\text{avg}}(T)=\frac{1}{N}\sum_{i=1}^{N}p_{i}(T) = 5.574\%.
\end{equation}
Table~\ref{tab:lhp_mc} reports the comparison at \(\rho=0.3\).
The equity tranche agrees to within 1.3~bp of tranche notional
(0.7\% relative), while mezzanine and senior tranches deviate by
44--69~bp. The mechanism is the heterogeneity of the pool
(\(p_{i}\in[0.859\%,\,47.778\%]\)): LHP collapses the pool into a
homogeneous Vasicek mass at \(p^{\text{avg}}\) and therefore
\emph{under-samples} the right tail of the true loss distribution.
This is the well-known pedagogical limitation of LHP
(\citealp[\S6.5]{mounfield2009}; \citealp{okane2008}); production
workflows move to bucketed LHP, recursive Hull--White or pure
Monte-Carlo for non-equity tranches. In our pipeline we retain LHP
only for equity sanity checks and for the historical dashboard of
Section~\ref{sec:dashboard}.

\begin{table}[H]
\centering
\caption{LHP vs Monte-Carlo expected tranche loss at
\(\rho=0.3\), valuation date 10~September~2021, MC=\num{10000}
paths.}
\label{tab:lhp_mc}
\setlength{\tabcolsep}{8pt}
\begin{tabular}{l c c c c c}
\toprule
\textbf{Tranche} & \textbf{Attach} & \textbf{Detach}
& \textbf{\(\ETL^{\text{LHP}}\)} & \textbf{\(\ETL^{\text{MC}}\)}
& \textbf{\(\Delta\) (bp of notional)}\\
\midrule
Equity        & 0\%  & 3\%  & 0.01847 & 0.01834 & $+1.3$\\
Mezzanine     & 3\%  & 7\%  & 0.01403 & 0.00909 & $+49.4$\\
Junior Senior & 7\%  & 10\% & 0.00746 & 0.00304 & $+44.2$\\
Senior        & 10\% & 15\% & 0.00911 & 0.00222 & $+68.8$\\
\bottomrule
\end{tabular}
\end{table}

%==============================================================================
\section{Valuation and Pricing Results}\label{sec:pricing}

This section presents the pricing output of the pipeline for the
valuation date \textbf{10~September~2021} under the Monte-Carlo engine
of Section~\ref{sec:mc}. All numbers are reported at a flat asset
correlation \(\rho = 0.30\), \num{10000} Monte-Carlo paths with
antithetic variates, recovery \(R = 40\%\) and flat continuously
compounded risk-free rate \(r = 2\%\).

\subsection{Baseline fair running spreads}
Table~\ref{tab:fair_spreads} reports the fair running spread
\(s^{\star}\), the expected tranche loss both as a fraction of
tranche notional and of deal notional, and the tranche width. The
capital-structure monotonicity
\(s^{\star}_{\mathrm{eq}} > s^{\star}_{\mathrm{mez}}
>s^{\star}_{\mathrm{jsr}}>s^{\star}_{\mathrm{sen}}\)
is strictly respected.

\begin{table}[H]
\centering
\caption{Baseline tranche pricing on 10~September~2021
(\(\rho=0.30\), \(T=5\)Y, \(N=651\) equal-weighted names,
Monte-Carlo \num{10000} paths, seed 42).}
\label{tab:fair_spreads}
\setlength{\tabcolsep}{6pt}
\begin{tabular}{l c c c c c c}
\toprule
\textbf{Tranche} & \textbf{Attach} & \textbf{Detach} & \textbf{Width}
& \textbf{\(\ETL(T)\) tranche} & \textbf{\(\ETL(T)\) notional}
& \textbf{Fair spread \(s^{\star}\)}\\
 & $A$ & $D$ & $D-A$ & (fraction) & (fraction) & (bp / year)\\
\midrule
Equity        & 0\%  & 3\%  & 3\% & 0.6114 & 0.01834 & 1\,952.6\\
Mezzanine     & 3\%  & 7\%  & 4\% & 0.2273 & 0.00909 & 503.3\\
Junior Senior & 7\%  & 10\% & 3\% & 0.1013 & 0.00304 & 202.7\\
Senior        & 10\% & 15\% & 5\% & 0.0445 & 0.00222 & 85.4\\
\bottomrule
\end{tabular}

\medskip
\begin{minipage}{0.95\textwidth}
\footnotesize \(\ETL(T)\) tranche is the expected loss normalised by
tranche notional \((D-A)\); \(\ETL(T)\) notional is the same quantity
as a fraction of deal (pool) notional. The equity tranche absorbs
61.1\% of its \(3\%\) notional in expectation --- reflecting the
first-loss nature of this layer --- and commands a running spread of
more than 1\,950~bp, about 23 times the senior spread.
\end{minipage}
\end{table}

\subsection{Clean price of a long-protection position}
For a long-protection position on tranche \([A,D]\) with contractual
running coupon \(s^{c}\) (typically the standardised 100~bp for
mezzanine/senior and 500~bp for equity on iTraxx/CDX), the present
value under the model is
\begin{equation}
\mathrm{PV}^{\text{long prot}}_{[A,D]}
\;=\; \bigl(s^{\star}-s^{c}\bigr)\cdot A_{\mathrm{trn}}\cdot(D-A),
\label{eq:long_prot_pv}
\end{equation}
where the risky annuity
\(A_{\mathrm{trn}}=V_{\text{prot}}/s^{\star}\,(D-A)\) is a by-product
of the fair-spread computation~\eqref{eq:fair_spread}. The first
factor \((s^{\star}-s^{c})\) in~\eqref{eq:long_prot_pv} measures the
richness or cheapness of the standard-coupon contract relative to
the model-fair level, and is the \emph{spread basis}.

\paragraph{Spread basis at the valuation date.}
Comparing the fair spreads of Table~\ref{tab:fair_spreads} to the
iTraxx-style standard coupons yields the following spread basis:
\(\Delta s_{\mathrm{eq}}= +1\,452.6\)~bp (equity vs 500~bp),
\(\Delta s_{\mathrm{mez}}= +403.3\)~bp,
\(\Delta s_{\mathrm{jsr}}= +102.7\)~bp and
\(\Delta s_{\mathrm{sen}}= -14.6\)~bp (vs 100~bp).  The senior
tranche is the only one richer than the standard coupon, consistent
with the post-2020 re-rating of systemic tail risk.

\subsection{Upfront quoting convention for the equity tranche}
Because the equity fair spread is far above the standardised
\(500\)~bp coupon, the market convention is to quote it as an
\emph{upfront} payment \(U_{\mathrm{eq}}\) (as a percentage of
tranche notional) plus the fixed running coupon:
\begin{equation}
U_{\mathrm{eq}}
\;=\;\bigl(s^{\star}_{\mathrm{eq}}-s^{c}\bigr)\cdot A_{\mathrm{eq,\,trn}}.
\label{eq:upfront_eq}
\end{equation}
The pipeline returns an equity \(\ETL(T)=61.14\%\) of tranche
notional, so a protection buyer at \(s^{c}=500\)~bp would pay an
indicative upfront of order
\(U_{\mathrm{eq}}\approx(1\,952.6-500)\,\mathrm{bp}\times A_{\mathrm{eq,\,trn}}
\approx 45\%\) of equity-tranche notional (using a representative
equity annuity of \(A_{\mathrm{eq,\,trn}}\approx 3.1\) years), a
level consistent with distressed equity-tranche marks observed in
practice.

\subsection{Mark-to-market of an existing position}\label{sec:mtm}
For a pre-existing long-protection position entered at contractual
spread \(s^{0}\), the mark-to-market at valuation date is simply
\begin{equation}
\mathrm{MTM}_{[A,D]}
\;=\; \bigl(s^{\star}_{t}-s^{0}\bigr)\cdot A_{\mathrm{trn},t}
\cdot\bigl(D-A\bigr),
\label{eq:mtm}
\end{equation}
so that positive values correspond to a widening of the
market-implied tranche spread relative to the trade entry.
Equation~\eqref{eq:mtm} is the hook that connects daily P\&L to the
CS01, Rho01 and JtD sensitivities of Section~\ref{sec:risk}.
For instance, a mezzanine long-protection leg opened in 2019 at
\(s^{0}=300\)~bp would mark at
\((503.3-300)\,\mathrm{bp}\times A_{\mathrm{trn}}\times 4\%\approx
0.26\%\) of deal notional on the 10~September~2021 snapshot.

\subsection{Convergence and Monte-Carlo noise}
The Monte-Carlo estimator of \(\ETL(T)\) converges at the standard
\(O(n^{-1/2})\) rate. Empirically, at the production setting
(\(n=\num{10000}\), antithetic variates, common random numbers across
tranches) the standard error on the fair spread is of order
\(0.5\)~bp for the equity tranche and below \(0.2\)~bp for the
senior tranche, small enough not to contaminate the
sensitivity analysis of Section~\ref{sec:risk}. Note, however,
that as documented above the antithetic
variate in isolation yields only marginal variance reduction
(\(\mathrm{VRF}\approx 0.98\)): the main sources of stability of
the estimator are common random numbers and the large pool size
(\(N=651\)).

\subsection{Worked example: the mezzanine tranche}
As a complete end-to-end numerical illustration, consider the
\emph{mezzanine} tranche \([3\%,\,7\%]\). The pipeline returns
\(\ETL_{5Y}^{\mathrm{mez}} = 0.2273\) of tranche notional (i.e.\
\(0.00909\) of deal notional) and a fair running spread
\(s^{\star}_{\mathrm{mez}} = 503.3\)~bp/yr. At the standard running
coupon \(s^{c}=100\)~bp, the clean PV of a freshly struck
long-protection position is, from \eqref{eq:long_prot_pv},
\begin{equation*}
\mathrm{PV}_{\mathrm{mez}} \;\approx\; (503.3-100)\,\mathrm{bp}
\cdot A_{\mathrm{trn}}\cdot (7\%-3\%)
\;\approx\; 4.03\%\times A_{\mathrm{trn}}\times 4\%
\;\approx\; 0.52\%\ \text{of deal notional},
\end{equation*}
taking \(A_{\mathrm{trn}}\!\approx\!3.2\) years as a representative
annuity for a 5-year mezzanine with cumulative loss \(\approx 23\%\).
A parallel \(+10\)~bp widening of all pool spreads shifts the fair
mezzanine spread by \(+96.4\)~bp (CS01 of \(+9.641\) per bp), confirming
the moderate first-order sensitivity of the tranche to systematic
credit risk.

%==============================================================================
\section{Base-Correlation Calibration and Bespoke Tranches}\label{sec:calib}

\subsection{Compound vs base correlation}
Given a market quote \(s^{\star,\mathrm{mkt}}_{[A_{j},D_{j}]}\), one
can either back out a \emph{compound correlation}
\(\rho^{\mathrm{cmp}}_{j}\) that re-prices the individual tranche, or
a \emph{base correlation}
\(\rho^{B}(D_{j})\) that re-prices the \emph{equity-gathered} tranche
\([0,D_{j}]\).  The latter, introduced by \citet{jpmorgan2004}, is
preferred because it is monotone in \(D\) and admits interpolation.
A bespoke tranche \([A,D]\) is then priced as
\begin{equation}
\ETL^{[A,D]}
= \ETL^{[0,D]}\!\bigl(\rho^{B}(D)\bigr)
- \ETL^{[0,A]}\!\bigl(\rho^{B}(A)\bigr).
\label{eq:bespoke_etl}
\end{equation}

\subsection{A pedagogical round-trip}
In the absence of broker quotes, we simulate a market smile along
the canonical iTraxx pattern, with base correlation rising with
detachment as a reflection of the systemic-risk premium:
\begin{equation}
\rho^{B}(D):\;
\{3\%\!\to\!20\%,\;7\%\!\to\!30\%,\;10\%\!\to\!37\%,\;15\%\!\to\!45\%\}.
\label{eq:bc_quote}
\end{equation}
For each detachment we generate the target
\(\ETL^{[0,D]}\) from the LHP-at-quote, and then recover the base
correlation curve through bisection on the same LHP objective.
Table~\ref{tab:bc_roundtrip} reports the outcome.  Because the
target ETLs are generated on the heterogeneous pool while the
objective function is evaluated on the LHP approximation, the
bisection saturates at the upper bracket \(\rho^{B}=99.9\%\) for
all four detachments: under the homogeneous LHP with
\(p^{\text{avg}}=5.574\%\), no level of correlation can reproduce
the fat-tailed \([0,D]\) ETL of the heterogeneous pool. This
illustrates exactly the heterogeneity caveat of
Section~\ref{sec:mc}: BC calibration against an LHP objective is
only sensible when the target itself is also LHP-generated (a true
round-trip) or when a bucketed LHP is used.  Robust production
calibration requires either (i) broker index-tranche quotes as
targets with a full MC pricer inside the objective or (ii) a
bucketed LHP surrogate with several homogeneous groups.

\begin{table}[H]
\centering
\caption{Base-correlation round-trip diagnostic,
10~September~2021. Target ETL is generated from the heterogeneous
pool at the market-BC quote; recovered BC uses the homogeneous LHP
objective.}
\label{tab:bc_roundtrip}
\begin{tabular}{c c c c}
\toprule
\textbf{Base tranche} & \textbf{Target ETL (bp)}
& \textbf{Market BC (input)} & \textbf{Recovered BC (LHP)}\\
\midrule
\([0,\,3\%]\)  & 225.2 & 20.00\% & 99.90\%\\
\([0,\,7\%]\)  & 329.8 & 30.00\% & 99.90\%\\
\([0,\,10\%]\) & 377.7 & 37.00\% & 99.90\%\\
\([0,\,15\%]\) & 450.7 & 45.00\% & 99.90\%\\
\bottomrule
\end{tabular}
\end{table}

\subsection{Bespoke pricing surface}
Once \(\rho^{B}(D)\) is known, any bespoke tranche \([A,D]\) with
\(A,D\in[0,15\%]\) is priced via \eqref{eq:bespoke_etl} and
\eqref{eq:fair_spread}, with piecewise-linear interpolation and flat
extrapolation of the base-correlation curve. The resulting
fair-spread surface \(s^{\star}(A,D-A)\) gives the desk a
fast-quoting tool over an arbitrary grid.

\begin{figure}[H]
    \centering
    % TODO: base correlation curve + bespoke surface
    % \includegraphics[width=0.45\textwidth]{fig_bc_curve.png}
    % \includegraphics[width=0.45\textwidth]{fig_bespoke_surface.png}
    \fbox{\parbox[c][5cm][c]{0.75\textwidth}{\centering\itshape
        Figure placeholder --- \textbf{(left)} Base-correlation curve
        \(\rho^{B}(D)\) and \textbf{(right)} bespoke fair-spread
        surface \(s^{\star}(A,D-A)\).}}
    \caption{Base-correlation curve and bespoke fair-spread surface.}
    \label{fig:bc_surface}
\end{figure}

%==============================================================================
\section{Risk Analytics}\label{sec:risk}

The pricing engine supports a full set of first-order sensitivities
(Greeks) through a bump-and-re-price methodology executed on
\emph{common random numbers} (identical paths) in order to control
Monte-Carlo noise.

\subsection{Tranche-level Greeks}
Table~\ref{tab:greeks} reports the full set of tranche-level
sensitivities at the 10~September~2021 snapshot, computed by bump
and re-price on common random numbers. All numerical predictions of
the economic intuition are verified:
(i) equity CS01 is the largest by a factor of 7 relative to the
senior tranche (first-loss leverage);
(ii) Rho01 is strictly \emph{negative} for the equity tranche
(\(-37.5\) bp per 1~point of \(\rho\)) and strictly \emph{positive}
for both senior layers, with the mezzanine almost Rho-neutral
(\(-0.73\)~bp) --- in line with the received wisdom that
\emph{equity is long correlation, senior is short correlation};
(iii) jump-to-default is largest for equity (one additional default
moves the equity fair spread by \(+145.5\)~bp) and collapses
towards zero for deep-out-of-the-money senior tranches (\(+1.5\)~bp).

\begin{table}[H]
\centering
\caption{Tranche-level Greeks at 10~September~2021,
\(\rho=0.30\), \(T=5\)Y, MC=\num{10000} paths.}
\label{tab:greeks}
\setlength{\tabcolsep}{8pt}
\begin{tabular}{l r r r r r}
\toprule
\textbf{Tranche} & \textbf{Fair spread} & \textbf{CS01}
& \textbf{Rho01} & \textbf{DV01} & \textbf{JtD}\\
 & (bp/yr) & (bp/+1bp) & (bp/+1\%) & (bp/+1bp rate) & (bp/1 default)\\
\midrule
Equity        & 1\,952.6 & $+21.926$ & $-37.502$ & $+0.026$ & $+145.54$\\
Mezzanine     &   503.3 & $+9.641$  & $-0.732$  & $-0.023$ & $+12.86$\\
Junior Senior &   202.7 & $+5.628$  & $+3.902$  & $-0.014$ & $+3.86$\\
Senior        &    85.4 & $+3.019$  & $+3.829$  & $-0.007$ & $+1.52$\\
\bottomrule
\end{tabular}
\end{table}

\subsection{Name-level CS01 and hedge concentration}
Under the LHP mean-field approximation, the CS01 of the
\(k\)-th name reads
\begin{equation}
\CS_{k}
 = \frac{\partial\ETL}{\partial p^{\text{avg}}}
   \cdot\frac{\partial p_{k}}{\partial s_{k}}\cdot\frac{1}{N}.
\label{eq:name_cs01}
\end{equation}
Empirically on our pool, the top-10 issuers by \(\CS_{k}\)
concentrate only \textbf{1.6\%} of the pool-level CS01 on the
equity tranche, and the top-50 issuers cumulatively capture
\textbf{8.0\%} (identical ratios on all other tranches because
under the LHP mean-field approximation the name-level contribution
is uniform in \(\partial p_k/\partial s_k/N\) and depends on the
tranche only through the common factor
\(\partial \ETL/\partial p^{\text{avg}}\)). This is a
\emph{benign} risk profile: the CS01 is highly diversified and no
small sub-portfolio of CDS can efficiently hedge the tranche
spread risk. A fully replicating CDS basket of several hundred
names is required --- a natural consequence of the portfolio
having \(N=651\) approximately equal-weighted issuers with
relatively tight PD dispersion around the pool mean.
Figure~\ref{fig:name_cs01} shows the rank-ordered top-20 CS01
contributors, led by names with the tightest 5Y PDs
(NOVNVX, AGSBB, XL, TOYOTA, NESNVX, \ldots with \(p_i\) near
\(0.9\%\)).

\begin{figure}[H]
    \centering
    % TODO: name-level CS01 bar chart
    % \includegraphics[width=0.8\textwidth]{fig_name_cs01.png}
    \fbox{\parbox[c][5cm][c]{0.75\textwidth}{\centering\itshape
        Figure placeholder --- Top-20 issuer contributions to the
        equity CS01. Hedge concentration: top-10 names capture roughly
        30--40\% of total equity CS01.}}
    \caption{Name-level CS01 and hedge concentration.}
    \label{fig:name_cs01}
\end{figure}

\subsection{Stress-testing matrix}
A regulator-grade deterministic stress grid is applied without
probabilistic weighting. The shocks follow the Basel III and Fed CCAR
templates:
\begin{itemize}[leftmargin=2em]
    \item Spread shock: parallel \(+50,+100,+200\) bp across the pool;
    \item Recovery shock: \(R\) reduced from 40\% to 25\% and to 10\%;
    \item Correlation shock: \(\rho\) increased by +20 and +50 points;
    \item Combined crisis shock: the superposition of spread \(+200\) bp,
          recovery \(25\%\) and \(\Delta\rho = +30\) pts.
\end{itemize}
The resulting heat-map of \(\Delta s^{\star}\) (Figure~\ref{fig:stress})
identifies the worst scenario for each tranche and supports
capital-planning and hedge-sizing decisions.

\begin{figure}[H]
    \centering
    % TODO: stress matrix heatmap
    % \includegraphics[width=0.8\textwidth]{fig_stress_matrix.png}
    \fbox{\parbox[c][5cm][c]{0.75\textwidth}{\centering\itshape
        Figure placeholder --- Heat-map of
        \(\Delta s^{\star}\) by (scenario, tranche).
        The combined 2008-like crisis is the worst scenario for all
        mezzanine and senior tranches.}}
    \caption{Deterministic stress-testing matrix.}
    \label{fig:stress}
\end{figure}

%==============================================================================
\section{Historical Dashboard: 2015--2021}\label{sec:dashboard}

To place the deal-date pricing in historical context we compute
81 monthly snapshots of the expected tranche loss under the
flat-hazard model and the LHP formula (\(n_{\text{gh}}=20\)) from
January~2015 through September~2021. The pool-average 5Y PD
oscillates between a floor of \(5.56\%\) and a peak of
\(\mathbf{11.84\%}\) reached on \textbf{1~April~2020}, at the COVID
credit shock. Two key empirical observations emerge.

\begin{itemize}[leftmargin=2em]
    \item \textbf{Tranche leverage during COVID.} The ratio of peak
          ETL (April~2020) to pre-COVID mean ETL (2019) quantifies
          the non-linear convexity of each tranche to the systemic
          shock: \(1.21\times\) for the equity, \(1.33\times\) for
          the mezzanine, \(1.39\times\) for the junior senior and
          \(\mathbf{1.45\times}\) for the senior tranche. As
          expected, the senior tranche exhibits the highest
          leverage --- senior tranches are short-correlation and
          short-gamma in the tail, so a pool-level PD shock
          translates into a disproportionate jump in tranche ETL
          (Figure~\ref{fig:hist_dashboard}).
    \item \textbf{Regime shifts.}
          The pre-COVID (2017--2019) and post-COVID recovery
          (2021) periods show a flat and low-volatility regime;
          Q1--Q2 2020 is a sharp spike lasting roughly one quarter
          before mean-reversion.
\end{itemize}

\begin{figure}[H]
    \centering
    % TODO: historical ETL dashboard
    % \includegraphics[width=0.9\textwidth]{fig_historical_dashboard.png}
    \fbox{\parbox[c][5cm][c]{0.80\textwidth}{\centering\itshape
        Figure placeholder --- Monthly tranche ETL, January~2015 --
        December~2021, by tranche; COVID shock shaded in red.
        Equity tranche scaled on the left axis, senior on the right.}}
    \caption{Historical expected tranche loss (LHP monthly snapshot).}
    \label{fig:hist_dashboard}
\end{figure}

%==============================================================================
\section{Validation Checklist}\label{sec:validation}

The pipeline ends with a single-page risk report and an automated
validation checklist. The latter contains four mandatory checks,
each producing a PASS/FAIL flag; the live outputs on the
10~September~2021 snapshot are summarised in
Table~\ref{tab:validation}.

\begin{table}[H]
\centering
\caption{Automated validation checklist
(10~September~2021).}
\label{tab:validation}
\begin{tabular}{c l l l}
\toprule
\textbf{\#} & \textbf{Check} & \textbf{Measured value} & \textbf{Status}\\
\midrule
1 & CDS par-repricing on bootstrap curve
  & $\max|\mathrm{NPV}|=8.33\!\times\!10^{-17}$, 100\% of 300 names
  & \textbf{PASS}\\
2 & LHP vs MC equity ETL relative error
  & $|\Delta|/\ETL^{\mathrm{MC}} = 0.7\%$ (target \(<20\%\))
  & \textbf{PASS}\\
3 & Capital-structure monotonicity
  & $1952.6 > 503.3 > 202.7 > 85.4$
  & \textbf{PASS}\\
4 & Correlation-sign test
  & $\mathrm{Rho01}_{\mathrm{eq}}=-37.5 <0$,
    $\mathrm{Rho01}_{\mathrm{sen}}=+3.83 >0$
  & \textbf{PASS}\\
\bottomrule
\end{tabular}
\end{table}

All four checks pass on the 10~September~2021 snapshot, confirming
that the pipeline is internally consistent and ready for sign-off by
model validation.

%==============================================================================
\section{Conclusion and Limitations}\label{sec:conclusion}

We have presented a complete and reproducible valuation and risk
pipeline for synthetic CDOs within the one-factor Gaussian copula
framework. The pipeline is internally consistent, benchmarked against
the LHP semi-analytic formula, calibrated to a base-correlation curve,
and enriched with a full suite of sensitivities and stress tests. The
historical dashboard highlights the non-linear leverage of senior
tranches through the COVID shock and quantitatively motivates the
inclusion of a combined crisis stress in the capital-planning grid.

Two limitations warrant explicit statement. First, the Gaussian
copula fails to reproduce the empirical smile of base correlations
and should be augmented, in production, with richer dependence
structures (double-\(t\), Marshall--Olkin, or a stochastic-recovery
extension; see \citealp{andersen2003,burtschell2007}). Second, the
LHP benchmark is valid only under homogeneity and displays a
material heterogeneity error on mezzanine and senior tranches for
real-world pools; bucketed LHP or the recursive Hull--White algorithm
\citep{hull2004} should be used when semi-analytic speed is required
for non-equity tranches.

\paragraph{Further work.}
Natural extensions include (i) a stochastic-recovery overlay
calibrated jointly with \(\rho^{B}\); (ii) a multi-period (forward
starting) tranche calibration to capture spread-term-structure
effects; (iii) a machine-learning surrogate of the LHP--MC difference
function to accelerate bespoke pricing and (iv) a scenario generator
based on PCA of the CDS time series to complement the deterministic
stress matrix.

%==============================================================================
\bigskip
\begin{thebibliography}{99}

\bibitem[Andersen and Sidenius(2003)]{andersen2003}
Andersen, L. and Sidenius, J. (2003).
\emph{Extensions to the Gaussian Copula: Random Recovery and Random
Factor Loadings}. Journal of Credit Risk, 1(1), 29--70.

\bibitem[Basel Committee on Banking Supervision(2006)]{bcbs2006}
Basel Committee on Banking Supervision (2006).
\emph{International Convergence of Capital Measurement and Capital
Standards --- A Revised Framework}. BCBS Publications.

\bibitem[Burtschell et al.(2007)]{burtschell2007}
Burtschell, X., Gregory, J. and Laurent, J.-P. (2007).
\emph{Beyond the Gaussian Copula: Stochastic and Local Correlation}.
Journal of Credit Risk, 3(1), 31--62.

\bibitem[Hull and White(2004)]{hull2004}
Hull, J. and White, A. (2004). \emph{Valuation of a CDO and an
n-th-to-Default CDS Without Monte Carlo Simulation}. Journal of
Derivatives, 12(2), 8--23.

\bibitem[JP Morgan(2004)]{jpmorgan2004}
JP Morgan (2004). \emph{Base Correlation}. Credit Derivatives
Research, JP Morgan.

\bibitem[Li(2000)]{li2000}
Li, D. X. (2000).
\emph{On Default Correlation: A Copula Function Approach}.
Journal of Fixed Income, 9(4), 43--54.

\bibitem[Mounfield(2009)]{mounfield2009}
Mounfield, C. C. (2009). \emph{Synthetic CDOs: Modelling, Valuation
and Risk Management}. Cambridge University Press.

\bibitem[O'Kane(2008)]{okane2008}
O'Kane, D. (2008). \emph{Modelling Single-Name and Multi-Name Credit
Derivatives}. Wiley Finance.

\bibitem[Vasicek(1987)]{vasicek1987}
Vasicek, O. (1987). \emph{Probability of Loss on Loan Portfolio}.
KMV Corporation.

\end{thebibliography}

%==============================================================================
\appendix
\section{Pseudocode: bootstrap and ETL engine}\label{app:pseudo}

\begin{algorithm}[H]
\caption{ISDA-consistent piecewise-constant hazard bootstrap.}
\begin{algorithmic}[1]
\Require Quotes $\{s_i^{T_k}\}$, discount curve $\DF(\cdot)$, recovery $R$
\Ensure  Piecewise-constant hazard $\{\lambda_i^{(k)}\}_{k=1}^{K}$
\For{each issuer $i$}
  \State $\lambda_i^{(0)} \gets 0$
  \For{$k = 1,\dots,K$}
    \State Solve $V_{\text{prot}}(T_k;\lambda_i^{(k)})-V_{\text{prem}}(T_k;\lambda_i^{(k)}) = 0$
    \Statex \hspace{1.2em} by bisection, holding $\lambda_i^{(1)},\dots,\lambda_i^{(k-1)}$ fixed.
  \EndFor
\EndFor
\State \Return $\{\lambda_i^{(k)}\}$
\end{algorithmic}
\end{algorithm}

\begin{algorithm}[H]
\caption{Monte-Carlo tranche ETL with antithetic variates.}
\begin{algorithmic}[1]
\Require  $\{p_i(T)\}$, $\rho$, tranche $[A,D]$, $n_{\text{paths}}$, seed
\Ensure   ETL curve $(\ETL_k)_{k=0}^{K}$
\State Draw $M_{1:n/2}\sim\Nor(0,1)$ and set $M_{n/2+1:n}=-M_{1:n/2}$
\State Draw $(\varepsilon_{1:n,i})\sim\Nor(0,1)$
\State $Y_{ji}\gets\sqrt{\rho}\,M_j+\sqrt{1-\rho}\,\varepsilon_{ji}$
\State $U_{ji}\gets\Phi(Y_{ji});\quad \tau_{ji}\gets -\log(1-U_{ji})/\lambda_i$
\For{each payment date $t_k$}
  \State $L_j^{\text{pool}}(t_k)\gets\sum_i \LGD_i\,\1_{\{\tau_{ji}\le t_k\}}/N$
  \State $L_j^{[A,D]}(t_k)\gets\min(\max(L_j^{\text{pool}}(t_k)-A,0),D-A)$
  \State $\ETL_k\gets\frac{1}{n}\sum_j L_j^{[A,D]}(t_k)$
\EndFor
\State \Return $(\ETL_k)$
\end{algorithmic}
\end{algorithm}

\end{document}
